\section*{الفصل \ref{ch:b2:elements-in-the-earth} --- العناصر في الأرض}

\begin{solution}{exo:b2:elements-in-the-earth:1}
يمثّل البوتاسيوم $2\times39.10/(2\times39.10+16.00) = 78.20/94.20 = 0.8301$ من كتلة
\ce{K2O}، ومنه $\num{5130}\times0.8301 = \qty{4259}{ppm}$ من K. ويمثّل التيتانيوم
$47.87/79.87 = 0.5993$ من \ce{TiO2}: $\num{27310}\times0.5993 = \qty{16370}{ppm}$ من Ti.
\end{solution}

\begin{solution}{exo:b2:elements-in-the-earth:2}
الألومنيوم محبّ للصخور وعصيّ على التبخّر: فقد احتفظت به الأرض السيليكاتية كله، بنسب قريبة
من الكوندريتية (1.16). والصوديوم محبّ للصخور لكنه متطاير (\qty{958}{K}): فقد
فُقد قبل التراكم أو أثناءه، وبقي منه 0.22. والنيكل يتكاثف مع المغنيسيوم
(\qty{1353}{K})، فلا يستطيع التطاير تفسير 0.077: إنه محبّ للحديد، وقد ذهب معظمه
إلى النواة. والإيريديوم عصيّ على التبخّر (\qty{1603}{K}) وشديد المحبّة للحديد:
0.0029، وكله تقريبًا في النواة.
\end{solution}

\begin{solution}{exo:b2:elements-in-the-earth:3}
$\mathrm{La}_N = 31/0.232 = 133.6$، $\mathrm{Yb}_N = 2.0/0.163 = 12.27$،
و$(\mathrm{La}/\mathrm{Yb})_N = 10.9$. القشرة غنية بكل \omterm{def:b2:elements-in-the-earth:ree}{العناصر الأرضية النادرة}،
وأغنى بالخفيفة منها بإحدى عشرة مرة منها بالثقيلة: فهي مكوّنة من مصاهير، وهذه
تأخذ \omterm{def:b2:elements-in-the-earth:ree}{العناصر الأرضية النادرة} الخفيفة غير المتوافقة بسهولة أكبر.
\end{solution}

\begin{solution}{exo:b2:elements-in-the-earth:4}
$\mathrm{Sm}_N = 4.7/0.145 = 32.4$، $\mathrm{Eu}_N = 1.0/0.0546 = 18.3$،
$\mathrm{Gd}_N = 4.0/0.198 = 20.2$؛ $\mathrm{Eu}^* = \sqrt{32.4\times20.2} = 25.6$
و$\mathrm{Eu}/\mathrm{Eu}^* = 18.3/25.6 = 0.72$، أي شذوذ سالب. فالبلاجيوكلاز، الذي
يستقبل \ce{Eu^2+}، تبلّر من المصاهير التي بنت القشرة العلوية وبقي
في الأعماق.
\end{solution}

\begin{solution}{exo:b2:elements-in-the-earth:5}
$\bar D_{\mathrm{La}} = 0.65\times0.058+0.18\times0.001+0.17\times0.001 = 0.0377 +
0.0002 + 0.0002 = 0.038$؛ $\bar D_{\mathrm{Lu}} = 0.3445+1.044+0.0306 = 1.42$. اللوتيتيوم
متوافق ($\bar D > 1$)، واللانثانوم شديد عدم التوافق. ويحمل الغارنت
$1.044/1.42 = 74$ في المئة من لوتيتيوم الصخر، مع أنه لا يمثّل إلا 18 في المئة من
الصخر.
\end{solution}

\begin{solution}{exo:b2:elements-in-the-earth:6}
للصوديوم، تبلغ Na/Mg في الأرض السيليكاتية $2670/\num{228000} = 0.01171$ وقيمتها في \omterm{def:b2:elements-in-the-earth:ci}{كوندريت CI} هي
$5010/\num{95870} = 0.05226$: فالنسبة 0.224. وللنيكل، $1960/\num{228000}
= 0.008596$ مقابل $\num{10640}/\num{95870} = 0.1110$: أي 0.0775. فالصوديوم
مستنزف بسبب تطايره (\qty{958}{K})؛ أمّا النيكل فيتكاثف فوق \qty{1300}{K} وهو
أشدّ استنزافًا من الصوديوم بثلاث مرات لأنه محبّ للحديد: فقد أخذت النواة معظمه.
\end{solution}

\begin{solution}{exo:b2:elements-in-the-earth:7}
La: $5.21\times0.058 = \qty{0.302}{ppm}$؛ Lu: $0.53\times0.53 = \qty{0.281}{ppm}$.
وقيم المعايرة تختزل في نسبة النسبتين:
$(\mathrm{La}/\mathrm{Lu})_{N,\mathrm{cpx}}/(\mathrm{La}/\mathrm{Lu})_{N,\mathrm{melt}} =
D_{\mathrm{La}}/D_{\mathrm{Lu}} = 0.058/0.53 = 0.11$: فنمط البلّورة أشدّ انحدارًا نحو
\omterm{def:b2:elements-in-the-earth:ree}{العناصر الأرضية النادرة} الثقيلة من نمط المصهور بتسع مرات.
\end{solution}

\begin{solution}{exo:b2:elements-in-the-earth:8}
القشرة: $(12/0.658)/(31/0.648) = 18.24/47.84 = 0.38$. البازلت الهاوايي:
$(18.10/0.658)/(15.20/0.648) = 27.51/23.46 = 1.17$. بازلت MORB: $7.96/8.04 = 0.99$. والقشرة
وحدها، بعجزها الشديد في النيوبيوم، تحمل بصمة مناطق الاندساس.
\end{solution}

\begin{solution}{exo:b2:elements-in-the-earth:9}
\omterm{def:b2:elements-in-the-earth:bse}{الوشاح البدائي}: $240/0.0203 = \num{11800}$. بازلت MORB: \mbox{\textdir TLT K $= 0.160\times0.830\times10^4 =
\qty{1328}{ppm}$}، و\mbox{\textdir TLT K/U $= \num{11200}$}. القشرة: \mbox{\textdir TLT K $= 2.80\times0.830\times10^4 =
\qty{23240}{ppm}$}، و\mbox{\textdir TLT K/U $= \num{8600}$}. يرتفع البوتاسيوم مئة ضعف من الوشاح إلى القشرة،
ومع ذلك تتغيّر K/U بأقل من الثلث: فالعنصران متقاربان في عدم توافقهما،
والانصهار لا يكاد يفصل بينهما. واليورانيوم عصيّ على التبخّر، فوفرته تُستنتج من
النموذج الكوندريتي؛ ثم تعطي K/U المقيسة على الحمم بوتاسيوم الكوكب، الذي لا يمكن،
لكونه متطايرًا، أن يؤخذ من الكوندريتات.
\end{solution}

\begin{solution}{exo:b2:elements-in-the-earth:10}
عند التوازن $C_{\mathrm{rock}} = \bar D\,C_{\mathrm{melt}}$ لكل عنصر (والكسر
الصغير من المصهور يترك تركيب الصخر دون تغيّر يُذكر)، ومنه
$(C_1/C_2)_{\mathrm{melt}} = (C_{1}/\bar D_1)/(C_{2}/\bar D_2)$ لتركيزي الصخر $C_1,
C_2$، والنسبة هي $\bar D_2/\bar D_1$. ومن أجل La وLu: $1.419/0.03805 = 37$. فمصهور
الوشاح الحامل للغارنت أغنى باللانثانوم نسبةً إلى اللوتيتيوم من مصدره بعامل 37: فقد
احتجز الغارنت \omterm{def:b2:elements-in-the-earth:ree}{العناصر الأرضية النادرة} الثقيلة، فصار نمط المصهور شديد الانحدار.
\end{solution}

\begin{solution}{exo:b2:elements-in-the-earth:11}
$C_{\mathrm{mix}} = fC_U + (1-f)C_M$ لكل عنصر: خطي في $f$. والنسبة
$\bigl(f\,\mathrm{Th}_U + (1-f)\mathrm{Th}_M\bigr)/\bigl(f\,\mathrm{Nb}_U +
(1-f)\mathrm{Nb}_M\bigr)$ خارج قسمة دالتين خطيتين، أي قطع زائد في $f$.
تبلغ \mbox{\textdir TLT Th/Nb $= 0.2$} حين $f(10.5-2.4) + (1-f)(0.404-1.048) = 0$، أي $8.1f = 0.644(1-f)$،
$f = 0.644/8.744 = 0.074$: فنحو 7 في المئة من القشرة ترفع Th/Nb من 0.077 إلى 0.2. ومن أجل
$f = 0.1$: \mbox{\textdir TLT Th $= 1.05 + 0.364 = 1.414$}، و\mbox{\textdir TLT Nb $= 1.2 + 4.716 = 5.916$}، و\mbox{\textdir TLT Th/Nb $= 0.24$}.
\end{solution}

\begin{solution}{exo:b2:elements-in-the-earth:12}
إزالة مكوّن لا يحوي أيًّا من العنصرين تترك كتلتيهما، ومن ثمّ نسبتهما،
دون تغيير. فإذا احتُفظ بكل الألومنيوم، فإن $(\mathrm{Al}/\mathrm{Mg}) = 1.16$ ضعف CI
تعني أن الأرض السيليكاتية تحوي $1/1.16 = 0.86$ من مغنيسيومها الكوندريتي: أي إن 14 في
المئة مفقودة. والسيليكون نسبةً إلى الألومنيوم يبلغ $0.83/1.16 = 0.72$ ضعف قيمة CI:
فيكون 28 في المئة من السيليكون الكوندريتي مفقودًا.
\end{solution}

\begin{solution}{pb:b2:elements-in-the-earth:1}
\textbf{1.} \mbox{\textdir TLT Fe/Mg $= \num{182800}/\num{95870} = 1.907$}؛ \mbox{\textdir TLT Ni/Mg $= 0.1110$}؛ \mbox{\textdir TLT Si/Mg $=
1.111$}؛ \mbox{\textdir TLT Al/Mg $= 0.0887$}.

\textbf{2.} $(\num{182800}/55.85)/(\num{106500}/28.09) = 3273/3791 = 0.863$ ذرّة حديد لكل
ذرّة سيليكون، في حدود 7 في المئة من قيمة الفوتوسفير 0.81: وهي مقارنة معقولة، لأن كوندريتات CI
تطابق الشمس في العناصر المكوّنة للصخور.

\textbf{3.} فقد مكوّن لا يحوي أيًّا من عنصري نسبة ما يترك
النسبة دون تغيير (\cref{exo:b2:elements-in-the-earth:12}): لا تتغيّر إلا التراكيز المطلقة.

\textbf{4.} المغنيسيوم محبّ للصخور ولا يدخل النواة؛ وهو عنصر رئيسي،
معيّن جيدًا في الكوندريتات وفي الوشاح كليهما. أمّا الحديد فهو العنصر المتقاسَم
بين النواة والوشاح: فلا يمكن أن يكون مرجعًا لنفسه.

\textbf{5.} العناصر التي تتكاثف معًا لم يفصل التطاير بينها في
السديم، فينبغي أن تحتفظ الأرض الكلية بنسبها. وتشذّ Al/Mg: فالألومنيوم
عصيّ على التبخّر (\qty{1653}{K})، وAl/Mg في الأرض السيليكاتية تبلغ 1.16 ضعف قيمة CI.

\textbf{6.} \mbox{\textdir TLT Fe $= 8.05\times55.85/71.85 = 6.26$} في المئة؛ \mbox{\textdir TLT Ni $= 0.25\times58.69/74.69 =
0.196$} في المئة.

\textbf{7.} \mbox{\textdir TLT Fe/Mg $= 6.26/22.8 = 0.2746$}؛ \mbox{\textdir TLT Ni/Mg $= 0.196/22.8 = 0.00860$}.

\textbf{8.} الحديد: $1 - 0.2746/1.907 = 0.856$، أي 85.6 في المئة مفقودة؛ والنيكل:
$1 - 0.00860/0.1110 = 0.923$، أي 92.3 في المئة.

\textbf{9.} الإيريديوم أشدّ محبّةً للحديد من النيكل، والنيكل أشدّ من الحديد:
فكلما فضّل العنصر الفلز أكثر، أخذ الفلز منه إلى النواة أكثر.

\textbf{10.} في السيليكات، في صورة \ce{Fe^2+} في الأوليفين والبيروكسين (8 في المئة من FeO):
فالحديد المؤكسَد إلى FeO يرتبط بالأكسجين كعنصر محبّ للصخور وبقي في
الوشاح؛ ولم يغص إلا الحديد الفلزي.

\textbf{11.} $(1-x)C_{\mathrm{Fe}} + xc_{\mathrm{Fe}} = r_{\mathrm{Fe}}(1-x)C_{\mathrm{Mg}}$؛
والقسمة على $(1-x)c_{\mathrm{Fe}}$ تعطي القضية.

\textbf{12.} $r_{\mathrm{Fe}}C_{\mathrm{Mg}} - C_{\mathrm{Fe}} = 1.907\times22.8 - 6.26 =
37.22$؛ $y = 37.22/85.5 = 0.4353$، $x = y/(1+y) = 0.303$.

\textbf{13.} $0.1110\times22.8 - 0.196 = 2.335$؛ $y = 2.335/5.2 = 0.4490$، $x = 0.310$.

\textbf{14.} $y = (85.5\times37.22 + 5.2\times2.335)/(85.5^2 + 5.2^2) =
3194.4/7337.3 = 0.4354$، $x = 0.303$.

\textbf{15.} تزن كل معادلة بمربع معاملها: $85.5^2/5.2^2 = 270$
مرة أكثر للحديد، العنصر الأوفر بكثير.

\textbf{16.} $37.22/2.335 = 15.9$، مقابل $85.5/5.2 = 16.4$ في النموذج: أي فرق
3 في المئة.

\textbf{17.} $0.303\times\num{5.972e24} = \qty{1.81e24}{kg}$.

\textbf{18.} $(32.5-30.3)/32.5 = 7$ في المئة أقل.

\textbf{19.} $y = 0.325/0.675 = 0.4815$، والحديد في النواة $= 37.22/0.4815 = 77.3$ في المئة،
والنيكل $= 2.335/0.4815 = 4.85$ في المئة: أي 18 في المئة من العناصر الخفيفة، نحو ضعف النسبة التي تسمح بها الكثافة، من 5 إلى
10 في المئة. فلا بدّ أن تكون Fe/Mg في الأرض أعلى قليلًا من قيمة CI،
أو أن يكون مغنيسيومها السيليكاتي أعلى من مغنيسيوم النموذج.

\textbf{20.} السيليكون والكبريت والكربون والأكسجين والهيدروجين. والسيليكون ناقص في
الأرض السيليكاتية: فنسبته المُعايَرة على المغنيسيوم، 0.83، أدنى من المستوى العصيّ.

\textbf{21.} $1.111\times22.8 - 21.0 = 4.331$؛ $4.331/0.4354 = 9.9$ في المئة Si في
النواة.

\textbf{22.} 9.9 مقابل 6 في المئة: قد تحوي النواة سيليكونًا أكثر مما في النموذج، أو
لم يتراكم جزء من السيليكون قطّ. فالسيليكون يتكاثف عند \qty{1310}{K}، أدنى قليلًا
من المغنيسيوم والحديد (1336 و\qty{1334}{K}): فربما بقي بعضه في الغاز.

\textbf{23.} النواة: $0.325\times0.5 = \qty{0.1625}{ppm}$ من كتلة الأرض؛ والأرض
السيليكاتية: $0.675\times0.001 = \qty{0.000675}{ppm}$. فالنواة تحوي
$0.1625/0.1632 = 99.6$ في المئة من الذهب.

\textbf{24.} تضع الكيمياء وحدها
\textbf{\qty[text-series-to-math=true]{30.3}{\percent}} من كتلة الأرض في نواتها،
في حدود 7 في المئة من قيمة علم الزلازل 32.5 في المئة.
\end{solution}
